\documentclass[a4paper,10pt,twoside]{article}
\usepackage{its_conf}
\begin{document}
\setcounter{section}{0}
\setcounter{figure}{0}
\setcounter{table}{0}
\setcounter{equation}{0}
\setcounter{secnumdepth}{1}
\setcounter{secnumdepth}{1}
% ======= Начало Ваших тезисов. Все, что выше - не изменять ! ===========

\topic{Пространственный отбор больших графов в редакторе диаграмм}

\information{Хаджинова~К.~А.}
{Белорусский государственный университет информатики и радиоэлектроники}
{Минск, Республика Беларусь}
{kseniyakhadzhynava@gmail.com}

% ========= Аннотация ==========
\annotation{Рассматривается пространственный отбор элементов при отображении больших графов в редакторе диаграмм. На графах до 50\,000 узлов сравнены перебор, равномерная сетка и quadtree: сетка ускоряет поиск элемента в 440 раз, а quadtree уступает ей по скорости и времени построения.}

\begin{multicols}{2} %НЕ УДАЛЯТЬ!!!

\section*{Введение}
В редакторах диаграмм модель процесса, данных или архитектуры представляется атрибутированным графом, узлы которого соответствуют элементам диаграммы. При росте числа узлов возрастает время поиска элемента под курсором и отбора видимых элементов, которое входит в бюджет кадра интерфейса.

Для больших графов применяют кластеризацию данных [1], уровни детализации, при которых плотные области графа агрегируются [2], и отображение только текущей области просмотра: так, в GraphMaps на каждом масштабе рисуются лишь элементы внутри видимой области [3], а в ASK-GraphView используется иерархическая кластеризация [4]; обзор методов визуализации графов дан в [5]. Стандартным способом вывода в веб-приложениях является SVG, например в библиотеке D3 [6]: каждый элемент становится узлом DOM, что делает массовую отрисовку дорогой.

Определение элемента под курсором и отбор видимых элементов выполняются при каждом движении мыши и каждом кадре анимации, поэтому их время входит в бюджет кадра, равный 16,7~мс при частоте 60 кадров в секунду.

Цель работы~-- описать конвейер визуализации больших графов и оценить вклад пространственного отбора. Ранее авторами была предложена конвейерная модель обработки диаграмм с единым графовым представлением [7]; здесь исследуется его отображение.

\section{Конвейер визуализации}
Диаграмма задается текстом на предметно-ориентированном языке, где каждая строка описывает узел или связь. Конвейер состоит из семи этапов, показанных на рис.~\ref{fig:arch}: разбор текста, построение внутреннего представления графа, пространственный индекс, планировщик видимости, поэтапный рендерер и слой взаимодействия.

\image{
[width=\columnwidth]{fig-architecture-ru.png}
\caption{Архитектура конвейера визуализации\label{fig:arch}}
}

Как видно из рис.~\ref{fig:arch}, действия пользователя возвращаются к планировщику видимости. Для поиска элемента под курсором и отбора видимых элементов сравниваются перебор, равномерная сетка и quadtree. Перебор проверяет все $n$ узлов, индекс ограничивает проверку узлами нужной области. Равномерная сетка делит плоскость на ячейки одинакового размера и хранит в каждой номера пересекающих ее узлов, поэтому поиск проверяет одну ячейку. Quadtree [8] рекурсивно делит область на четыре части, пока в листе не окажется не больше заданного числа узлов. Эффективность отсечения невидимых элементов оценивается коэффициентом по формуле~(1):
$$
C_{cull} = 1 - N_{drawn}/N_{total},\eqno(1)
$$
где $N_{drawn}$ и $N_{total}$~-- число отображаемых и всех элементов.

Для поэтапного вывода узлу присваивается приоритет по формуле~(2):
$$
p(v) = 0{,}5\,deg(v) + 0{,}3\,level(v) + 0{,}2\,sel(v),\eqno(2)
$$
где $deg(v)$ и $level(v)$~-- нормализованные степень и уровень иерархии, $sel(v)$~-- признак выделения. Группы узлов с низким приоритетом (2) заменяются кластером с числом скрытых элементов.

\section{Методика эксперимента}
Измерения выполнены на отдельном стенде (Node.js 24, движок V8; Intel Core i5-1035G1, 24~ГБ ОЗУ, Windows~11). Генератор строит графы из 500, 2\,000, 10\,000 и 50\,000 узлов трех типов: случайных, иерархических и с узлами-хабами; начальное значение фиксировано. Измеряются время построения индекса, время поиска элемента под курсором (1\,000 точек) и время отбора узлов в окне $1600\times900$ (200 окон). Каждое измерение выполняется в отдельном процессе, 11 запусков, первый отбрасывается; приводятся медианы, для трех типов графов~-- медиана по типам. Результаты индексов сверены с перебором на 3\,000 узлов. Перед измерением выполняется прогрев (200 запросов поиска и 20 запросов отбора), время запроса измеряется функцией performance.now. Параметры индексов: ячейка сетки равна 120 условным единицам, лист quadtree вмещает до 8 узлов, глубина дерева не превышает 12.

\section{Результаты}
Результаты поиска элемента для 50\,000 узлов приведены в табл.~\ref{tab:idx}; индекс строится один раз при загрузке, поэтому указано и время его построения.

\myTable{Поиск элемента, 50\,000 узлов\label{tab:idx}}{{\columnwidth}{|>{\raggedright\arraybackslash}p{2.9cm}|X|X|X|} \hline
Метод & Индекс, мс & Поиск, мкс & 95\,\%, мкс \\ \hline
Перебор, объекты & 0 & 1446 & 4952 \\ \hline
Перебор, массивы & 0 & 573 & 2086 \\ \hline
Сетка & 69 & 1,3 & 4,4 \\ \hline
Quadtree & 644 & 5,0 & 18,8 \\ \hline
}

Равномерная сетка ускоряет поиск в 440 раз по сравнению с перебором по массивам, quadtree~-- в 115 раз (рис.~\ref{fig:hit}). Сетка быстрее quadtree и на графах с узлами-хабами (1,3 и 4,1~мкс против 621~мкс при переборе), а строится в 9 раз быстрее.

\image{
[width=0.8\columnwidth]{fig-hit-ru.png}
\caption{Время поиска элемента под курсором\label{fig:hit}}
}

На рис.~\ref{fig:hit} при росте числа узлов в 100 раз время поиска перебором по массивам возрастает в 155 раз, сеткой~-- в 2,4 раза, quadtree~-- в 3,3 раза. В окне $1600\times900$ остается около 190 узлов независимо от размера графа, поэтому по формуле (1) коэффициент $C_{cull}$ растет с 0,61 до 0,996 (табл.~\ref{tab:cull}).

\myTable{Отбор видимых узлов, мкс\label{tab:cull}}{{\columnwidth}{|X|X|X|X|X|} \hline
Узлов & $C_{cull}$ & Перебор & Сетка & Quadtree \\ \hline
500 & 0,61 & 7,0 & 7,1 & 18,8 \\ \hline
2\,000 & 0,90 & 38,8 & 9,3 & 24,3 \\ \hline
10\,000 & 0,98 & 156 & 20,4 & 36,2 \\ \hline
50\,000 & 0,996 & 775 & 49,9 & 55,4 \\ \hline
}

Время отбора на рис.~\ref{fig:cull} при росте числа узлов в 100 раз у перебора возрастает в 110 раз (с 7,0 до 775~мкс), у сетки~-- в 7 раз, у quadtree~-- в 3 раза.

\image{
[width=0.8\columnwidth]{fig-cull-ru.png}
\caption{Время отбора видимых узлов\label{fig:cull}}
}

На 500 узлах индекс не дает выигрыша, а начиная с 2\,000 узлов сетка выполняет отбор быстрее перебора. Даже при 50\,000 узлов перебор занимает меньше 1~мс и укладывается в бюджет кадра 16,7~мс, поэтому узким местом ожидается отрисовка; проверка для SVG и Canvas является следующим этапом. Из данных следует практическое правило: до 500 узлов достаточно перебора, начиная с 2\,000 узлов следует применять сетку, а quadtree оправдан лишь при сильно неравномерной плотности, которая не наблюдалась даже на графах с узлами-хабами.

\section*{Заключение}
Предложен конвейер визуализации больших графов и измерено влияние пространственного индекса. При близкой к равномерной плотности элементов предпочтительнее сетка: она в 3,8 раза быстрее при поиске и в 9 раз быстрее строится, чем quadtree. Ограничение работы~-- измерения выполнены в среде Node.js без отрисовки в браузере. Дальнейшая работа~-- измерение отрисовки SVG и Canvas и поэтапный вывод.

\ListReferences{
\item Козарь,~Р.~В. Методы кластеризации данных при распознавании эндоскопических снимков в задачах компьютерной медицинской диагностики~/ Р.~В.~Козарь, Н.~С.~Конойко, А.~А.~Навроцкий~// Доклады БГУИР.~-- 2023.~-- Т.~21, №~1.~-- С.~94--97.~-- DOI: 10.35596/1729-7648-2023-21-1-94-97.
\item Zinsmaier,~M. Interactive level-of-detail rendering of large graphs~/ M.~Zinsmaier, U.~Brandes, O.~Deussen, H.~Strobelt~// IEEE Trans. Vis. Comput. Graph.~-- 2012.~-- Vol.~18, №~12.~-- P.~2486--2495.
\item GraphMaps: browsing large graphs as interactive maps~/ L.~Nachmanson [et al.]~// arXiv.~-- 2015.~-- arXiv:1506.06745.
\item Abello,~J. ASK-GraphView: a large scale graph visualization system~/ J.~Abello, F.~van Ham, N.~Krishnan~// IEEE Trans. Vis. Comput. Graph.~-- 2006.~-- Vol.~12, №~5.~-- P.~669--676.
\item Herman,~I. Graph visualization and navigation in information visualization: a survey~/ I.~Herman, G.~Melancon, M.~S.~Marshall~// IEEE Trans. Vis. Comput. Graph.~-- 2000.~-- Vol.~6, №~1.~-- P.~24--43.
\item Bostock,~M. D3: data-driven documents~/ M.~Bostock, V.~Ogievetsky, J.~Heer~// IEEE Trans. Vis. Comput. Graph.~-- 2011.~-- Vol.~17, №~12.~-- P.~2301--2309.
\item Ломако,~А.~В. Автоматизация преобразования изображений диаграмм в программный код~/ А.~В.~Ломако, К.~А.~Хаджинова~// Технологии передачи и обработки информации.~-- 2026.
\item Finkel,~R.~A. Quad trees: a data structure for retrieval on composite keys~/ R.~A.~Finkel, J.~L.~Bentley~// Acta Inform.~-- 1974.~-- Vol.~4, №~1.~-- P.~1--9.
}

\end{multicols} % НЕ УДАЛЯТЬ!

% ===== Конец материалов. Все, что ниже, не изменять !===========
\end{document}
